
\section{Experimentos}

El diseño experimental de este trabajo se ha estructurado para evaluar el impacto de la complejidad arquitectónica y la transferencia de conocimiento 
en la tarea de clasificación de artrosis de rodilla. Para ello, se han entrenado y comparado cinco modelos representativos bajo condiciones homogéneas.

\textbf{Comparación de modelos con distinta complejidad.}
Se han considerado cinco arquitecturas con creciente capacidad de representación: tres modelos construidos desde cero (CNN simple, CNN pequeña y CNN profunda) 
y dos modelos preentrenados (EfficientNetB0 y ResNet50), ajustados mediante \textit{transfer learning}. Esta selección permite analizar la relación entre la 
expresividad del modelo, el coste computacional y la capacidad de generalización, tanto dentro del dominio (OAI) como fuera de él (radiografías de gatos).

\textbf{Configuración común de entrenamiento.}
Todos los modelos se han entrenado con los mismos hiperparámetros base: optimizador Adam, tasa de aprendizaje inicial de $10^{-3}$, función de pérdida
\texttt{CrossEntropyLoss}, tamaño de lote de 32 y un máximo de 50 épocas. Se ha aplicado \textit{early stopping} con paciencia de 10 épocas sin mejora 
en el conjunto de validación, y un \textit{scheduler} de tasa de aprendizaje con factor de reducción de 0.5 tras 5 épocas sin mejora.

En el caso de los modelos preentrenados, se han evaluado dos variantes: (i) congelando inicialmente el \textit{backbone} y afinando solo las capas finales; 
y (ii) realizando un ajuste fino completo. En ambos casos se ha empezado con una tasa de aprendizaje más baja ($10^{-4}$) debido a la mayor complejidad de
los modelos y la posibilidad de sobreajuste. Esta estrategia busca maximizar la transferencia de conocimiento y minimizar el riesgo de inestabilidad durante el
entrenamiento.

\textbf{División de datos y normalización.}
Todos los entrenamientos se han realizado utilizando las particiones oficiales del dataset OAI (\texttt{train}, \texttt{validation}, \texttt{test}), 
manteniendo su integridad para asegurar la replicabilidad. Las imágenes fueron redimensionadas a 224x224 píxeles y normalizadas para su compatibilidad 
con modelos preentrenados.

\textbf{Evaluación en dominio externo: dataset felino.}
Como medida de validación externa, se ha utilizado un conjunto adicional de imágenes de radiografías de gatos. Durante el entrenamiento, este conjunto 
se ha empleado como evaluación secundaria (no influenciada por el \textit{early stopping} ni el \textit{scheduler}) para estimar la capacidad de generalización 
de los modelos a un dominio visual distinto pero clínicamente análogo. Para cada modelo se ha conservado la mejor versión según la validación humana 
y otra según el rendimiento en gatos.

\textbf{Métricas de evaluación.}
El rendimiento de cada arquitectura se ha evaluado mediante varias métricas: precisión global (\textit{accuracy}), F1-score macro, matriz de confusión 
y error absoluto medio (MAE) para formulaciones ordinales. Estas métricas se han aplicado tanto al conjunto de test del dataset OAI como al conjunto 
felino, permitiendo un análisis cruzado entre dominios. Para la comparación de modelos se ha utilizado la precisión global como métrica principal.

\textbf{Técnicas de regularización y aumento de datos.}
Los modelos diseñados desde cero incluyen mecanismos de regularización como \texttt{Dropout} y \texttt{BatchNormalization}, cruciales para mejorar la 
estabilidad del entrenamiento y mitigar el sobreajuste. Además, se ha implementado \textit{data augmentation} en tiempo real: rotaciones aleatorias, 
reflejos verticales y horizontales, y variaciones de brillo y contraste. Estas transformaciones buscan enriquecer la variabilidad del conjunto de entrenamiento 
y mejorar la robustez del modelo.

\textbf{Resultados obtenidos}
Este apartado presenta de forma detallada los resultados experimentales obtenidos a lo largo del desarrollo del presente Trabajo de Fin de Grado. La evaluación de los modelos diseñados se ha centrado en analizar el impacto de la arquitectura, la formulación del problema (clasificación vs. regresión) y la partición de tareas (multiclase, binaria y ordinal) sobre la capacidad del sistema para predecir de forma automática el grado de artrosis de rodilla a partir de radiografías del conjunto de datos OAI.

En línea con los estudios revisados en el capítulo anterior, los experimentos se han concebido de forma progresiva, comenzando por modelos simples construidos desde cero y avanzando hacia arquitecturas más complejas y robustas. Esta progresión arquitectónica permite establecer una base comparativa sólida y observar cómo se incrementa la capacidad de representación de los modelos a medida que se introducen nuevas técnicas de normalización, regularización y profundidad.

\textbf{Evaluación de modelos simples construidos desde cero.}
El primer bloque de resultados corresponde a las tres CNN personalizadas (simple, mediana y profunda), entrenadas desde cero. El propósito inicial de estos modelos no es alcanzar una precisión máxima, sino establecer una línea base interpretativa que permita entender qué nivel de rendimiento puede obtenerse sin recurrir al aprendizaje por transferencia ni a arquitecturas sofisticadas. Así, la \emph{Simple CNN}, con una única capa convolucional y una densidad de parámetros mínima, permite observar la capacidad de generalización más básica del sistema.

En las tareas multiclase (KL 0 a 4), el modelo simple logra una precisión del 42\% y un F1-score del 30\%, lo que evidencia sus limitaciones para discriminar clases ordinales cercanas. No obstante, en tareas binarias extremas como KL 0 vs KL 4, la precisión asciende al 97.18\%, con un F1-score de 97.02\%, revelando que incluso modelos muy básicos pueden captar diferencias marcadas cuando las clases son visualmente muy dispares. Este comportamiento sugiere que la dificultad no reside únicamente en el modelo, sino en la naturaleza ambigua de los casos intermedios (KL1, KL2), en línea con lo reportado por Rafique et al. (2024).

La \emph{CNN mediana}, con dos capas convolucionales y una capa densa intermedia, introduce una mejora sustancial al permitir la detección de patrones jerárquicos más complejos. Su rendimiento en clasificación multiclase asciende al 46.85\%, con un F1-score del 39.71\%. La mejora más significativa se observa en tareas como 0 vs 3 y 0 vs 4, donde supera el 84\% y el 97\% de precisión, respectivamente. Estos resultados reflejan que con un incremento moderado en la profundidad y la capacidad de abstracción del modelo, se logran mejoras tangibles en la discriminación de clases intermedias.

Por su parte, la \emph{CNN profunda}, estructurada en tres bloques con \texttt{BatchNormalization} y \texttt{Dropout}, consigue consolidarse como la mejor red entre las diseñadas ad-hoc. En clasificación multiclase alcanza una precisión del 58.47\% y un F1-score del 52.03\%, lo que representa una mejora absoluta de más de 16 puntos respecto al modelo simple. En tareas binarias, destaca especialmente en las comparaciones intermedias (0 vs 2, 0 vs 3), donde su rendimiento supera ampliamente el 85\%, y alcanza el 99.72\% en la comparación 0 vs 4, con un F1-score idéntico. Estos resultados evidencian el papel crucial que desempeñan las técnicas de regularización en la mejora de la generalización.

\textbf{Resultados en formulación ordinal (regresión).}
Además de la clasificación categórica, se ha explorado una formulación regresiva del problema, utilizando las mismas arquitecturas pero adaptando la función de pérdida a un error cuadrático medio (MSE). Esta aproximación tiene como objetivo capturar la continuidad inherente al sistema de gradación Kellgren-Lawrence, penalizando más severamente los errores proporcionales a la distancia ordinal entre la predicción y la etiqueta real.

En esta configuración, los modelos muestran una tendencia coherente con lo observado en clasificación: la CNN profunda obtiene el mejor MAE (0.651), seguida por la CNN mediana (0.771) y la simple (0.897). De manera interesante, se observa que la formulación regresiva permite obtener mejores resultados en tareas reagrupadas (0-3-4), donde la CNN profunda alcanza un 84.34\% de precisión y un F1-score de 77.17\%, siendo este un resultado comparable al de estudios como Fei et al. (2024), quienes defienden también el enfoque continuo como una vía prometedora en tareas clínicas.

En esta segunda parte se presentan los resultados obtenidos al aplicar técnicas de \textit{transfer learning} sobre dos arquitecturas profundas ampliamente utilizadas en visión por computador: EfficientNetB0 y ResNet50. A diferencia de los modelos diseñados desde cero, estas redes aprovechan conocimiento previamente aprendido sobre grandes bases de datos como ImageNet, lo que permite partir de representaciones visuales genéricas y afinarlas para la tarea específica de clasificación de la artrosis.

\textbf{Resultados con EfficientNetB0.}
Este modelo ha destacado especialmente por su eficiencia y capacidad de generalización con un número relativamente reducido de parámetros. En su configuración óptima (ajuste fino completo), alcanza una precisión del 68.64\% en la tarea multiclase (KL 0-4), con un F1-score del 64.82\%. Estas cifras superan ampliamente las obtenidas por las CNN propias, con una mejora de más de 10 puntos porcentuales respecto al mejor modelo desde cero.

En las tareas binarias, EfficientNetB0 obtiene resultados notablemente altos, con una precisión del 94.07\% en la tarea 0-3-4 y del 97.70\% en 0 vs 3, consolidando su robustez en escenarios de clasificación clínica práctica. La formulación ordinal mediante regresión también resulta efectiva, alcanzando un MAE de solo 0.463 y una precisión del 65.13\%, lo que respalda la hipótesis de que estas arquitecturas permiten capturar mejor la continuidad inherente al grado de artrosis.

\textbf{Evaluación cruzada en el dominio felino.}
Uno de los objetivos específicos de este TFG ha sido evaluar la capacidad de generalización de los modelos desarrollados a un dominio anatómicamente distinto: el conjunto de imágenes radiográficas de gatos. Aunque este conjunto es de menor tamaño y presenta mayor variabilidad en las condiciones de adquisición, se ha mantenido como evaluación externa sin intervenir en los procesos de entrenamiento o validación.

Los resultados obtenidos muestran que los modelos preentrenados (especialmente EfficientNetB0) conservan un nivel razonable de precisión en este dominio externo, validando así la aplicabilidad del aprendizaje por transferencia inter-especie. Esta observación abre nuevas líneas de investigación en contextos veterinarios, donde los conjuntos de datos suelen ser mucho más escasos.

\textbf{Comparativa global y análisis.}
En la Figura~\ref{fig:comparativa_modelos_accuracy} se presenta una comparativa gráfica de la precisión multiclase (KL 0-4) alcanzada por los distintos modelos, permitiendo visualizar de forma directa la evolución del rendimiento conforme aumenta la complejidad de la arquitectura. Del mismo modo, la Tabla~\ref{tab:resumen_global_modelos} recoge de forma sintética las principales métricas (accuracy, F1, MAE) para cada configuración relevante.


Como puede observarse, la precisión y el F1-score mejoran de forma sostenida con el aumento de complejidad, y los modelos preentrenados presentan una clara ventaja en formulaciones multiclase. En tareas binarias extremas, sin embargo, incluso las CNN más simples alcanzan niveles notables, lo que refuerza la idea de que el grado de separación entre clases influye significativamente en la dificultad del problema.

\section{Conclusiones}
Los resultados obtenidos no solo confirman la viabilidad del uso de redes neuronales profundas para la clasificación automática de la artrosis de rodilla, sino que también validan el enfoque de este TFG al estructurar la evaluación de forma progresiva y comprensiva. La introducción de una evaluación interdominio basada en radiografías de gatos, aunque secundaria en este capítulo, ha servido para comprobar la robustez de los modelos desarrollados, lo que constituye una de las principales contribuciones del trabajo.

En los capítulos siguientes se analizarán más a fondo las interpretaciones clínicas de los resultados, así como la visualización de activaciones mediante Grad-CAM, con el objetivo de avanzar hacia sistemas de ayuda a la decisión médica más transparentes y explicables.

\textbf{Resultados simple CNN:}

\begin{table}[h]
    \centering
    \caption{Resultados Simple CNN en clasificación y regresión}
    \label{tab:simple_cnn_results_improved}
    \begin{tabular}{@{} l l l c c c c c c @{}} 
      \toprule
      \textbf{Dataset} & \textbf{Tarea} & \textbf{Métrica} 
        & \multicolumn{6}{c}{\textbf{Particiones}} \\
      \cmidrule(lr){4-9}
      & & & 0--4 & 0-3-4 & 0 vs 1 & 0 vs 2 & 0 vs 3 & 0 vs 4 \\
      \midrule
      % Clasificación OAI
      \multirow{2}{*}{OAI}
        & Clasif. & Accuracy   & 42\%   & 85.35\%   & 68.19\% &  65.37\%   &  80.88\%   &  97.18\%   \\
        & Clasif. & F1-score    & 30\%  & 80.35\%  & 55.29\%  &  62.84\%   &  77.98\%   &  97.02\%   \\
      \addlinespace
      % Regresión OAI
      \multirow{3}{*}{OAI}
        & Regr.   & Accuracy   & 34.87\%   & 84.38\% &  --      &  --   &  --   &  --   \\
        & Regr.   & F1-score    & 34.65\%  & 78.16\%  &  --      &  --   &  --   &  --   \\
        & Regr.   & MAE         & 0.897  & 0.291  &  --      &  --   &  --   &  --   \\
      \addlinespace
    \end{tabular}
\end{table}
  


\textbf{Resultados CNN mediana:}

\begin{table}[h]
    \centering
    \caption{Resultados CNN mediana en clasificación y regresión}
    \label{tab:cnn_mediana_results}
    \begin{tabular}{@{} l l l c c c c c c @{}} 
      \toprule
      \textbf{Dataset} & \textbf{Tarea} & \textbf{Métrica} 
        & \multicolumn{6}{c}{\textbf{Particiones}} \\
      \cmidrule(lr){4-9}
      & & & 0--4 & 0-3-4 & 0 vs 1 & 0 vs 2 & 0 vs 3 & 0 vs 4 \\
      \midrule
      % Clasificación OAI
      \multirow{2}{*}{OAI}
        & Clasif. & Accuracy   & 46.85\% & 85.35\% & 68.19\% & 62.41\% & 84.10\% & 97.46\% \\
        & Clasif. & F1-score   & 39.71\% & 80.35\% & 55.29\% & 54.41\% & 82.16\% & 97.35\% \\
      \addlinespace
      % Regresión OAI
      \multirow{3}{*}{OAI}
        & Regr.   & Accuracy   & 34.87\% & 84.02\% & -- & -- & -- & -- \\
        & Regr.   & F1-score   & 34.65\% & 82.71\% & -- & -- & -- & -- \\
        & Regr.   & MAE        & 0.771 & 0.237 & -- & -- & -- & -- \\
      \addlinespace
    \end{tabular}
\end{table}

\textbf{Resultados CNN grende:}
 
\begin{table}[h]
    \centering
    \caption{Resultados CNN grende en clasificación y regresión}
    \label{tab:cnn_grande_results}
    \begin{tabular}{@{} l l l c c c c c c @{}} 
      \toprule
      \textbf{Dataset} & \textbf{Tarea} & \textbf{Métrica} 
        & \multicolumn{6}{c}{\textbf{Particiones}} \\
      \cmidrule(lr){4-9}
      & & & 0--4 & 0-3-4 & 0 vs 1 & 0 vs 2 & 0 vs 3 & 0 vs 4 \\
      \midrule
      % Clasificación OAI
      \multirow{2}{*}{OAI}
        & Clasif. & Accuracy   & 58.47\% & 89.35\% & 68.40\% & 74.26\% & 86.64\% & 99.72\% \\
        & Clasif. & F1-score   & 52.03\% & 88.50\% & 55.78\% & 72.51\% & 86.14\% & 99.72\% \\
      \addlinespace
      % Regresión OAI
      \multirow{3}{*}{OAI}
        & Regr.   & Accuracy   & 43.95\% & 84.34\% & -- & -- & -- & -- \\
        & Regr.   & F1-score   & 46.22\% & 77.17\% & -- & -- & -- & -- \\
        & Regr.   & MAE        & 0.651 & 0.171 & -- & -- & -- & -- \\
      \addlinespace
    \end{tabular}
\end{table}

\textbf{Resultados EfficientNetB0:}
\begin{table}[h]
    \centering
    \caption{Resultados EfficientNetB0 en clasificación y regresión}
    \label{tab:cnn_grande_results}
    \begin{tabular}{@{} l l l c c c c c c @{}} 
      \toprule
      \textbf{Dataset} & \textbf{Tarea} & \textbf{Métrica} 
        & \multicolumn{6}{c}{\textbf{Particiones}} \\
      \cmidrule(lr){4-9}
      & & & 0--4 & 0-3-4 & 0 vs 1 & 0 vs 2 & 0 vs 3 & 0 vs 4 \\
      \midrule
      \multirow{2}{*}{OAI}
        & Clasif. & Accuracy   & 68.64\% & 94.07\% & 74.01\% & 87.22\% & 97.70\% & 99.72\% \\
        & Clasif. & F1-score   & 64.82\% & 93.84\% & 71.78\% & 87.77\% & 97.67\% & 99.72\% \\
      \addlinespace
      \multirow{3}{*}{OAI}
        & Regr.   & Accuracy   & 65.13\% & 93.83\% & -- & -- & -- & -- \\
        & Regr.   & F1-score   & 65.84\% & 93.67\% & -- & -- & -- & -- \\
        & Regr.   & MAE        & 0.463   & 0.102   & -- & -- & -- & -- \\
      \addlinespace
    \end{tabular}
\end{table}

\begin{table}[h]
    \centering
    \caption{Resultados ResNet151 en clasificación y regresión}
    \label{tab:cnn_grande_results}
    \begin{tabular}{@{} l l l c c c c c c @{}} 
      \toprule
      \textbf{Dataset} & \textbf{Tarea} & \textbf{Métrica} 
        & \multicolumn{6}{c}{\textbf{Particiones}} \\
      \cmidrule(lr){4-9}
      & & & 0--4 & 0-3-4 & 0 vs 1 & 0 vs 2 & 0 vs 3 & 0 vs 4 \\
      \midrule
      \multirow{2}{*}{OAI}
        & Clasif. & Accuracy   & 67.07\% & 95.16\% & 68.61\% & 86.67\% & 96.31\% & 100\% \\
        &         & F1-score   & 67.19\% & 94.89\% & 65.26\% & 86.55\% & 96.26\% & 100\% \\
      \addlinespace
      \multirow{3}{*}{OAI}
        & Regr.   & Accuracy   & 59.56\% & 82.93\% &  &  &  &  \\
        &         & F1-score   & 61.40\% & 79.17\% &  &  &  &  \\
        &         & MAE        & 0.515   & 0.342 &  &  &  &  \\
      \addlinespace
    \end{tabular}
\end{table}
